\section{Interpretive Framework and Next Steps}
\label{sec:discussion}

Although results are pending, the implemented design clarifies what each
eventual observation can and cannot mean.

\paragraph{Normalisation trades coverage for warranted claims.}
A low exact-cell yield would not by itself show poor annotation.  The closed
representation requires six scalar properties, whereas an EGP descriptor may
intentionally state only one contrast.  Retaining a partial mapping can be the
correct outcome.  RQ1 must therefore combine attrition rates with a manually
reviewed taxonomy: unsupported completion, missed evidence, wrong scope,
branching error, schema error, and genuine source underspecification.  The
right target is calibrated conservatism, not maximum downstream cell count.

\paragraph{Structural selection establishes sample support, not cognitive
truth.}
If the selector retains a candidate, the result means that its activation
pattern contributes to the declared development obligations after equivalence
and scope checks.  It does not show that learners acquire that candidate as an
indivisible skill.  Conversely, a rejected candidate may be meaningful outside
the small inventory but unsupported or indistinguishable within it.  The
selection trace, equivalence classes, and holdout activation patterns are at
least as important as the final KC count.

\paragraph{Coverage is part of predictive performance.}
Exact-cell KCs offer high structural specificity but cannot transfer learner
evidence to a new cell unless that cell was admitted in advance.  Reusable
factors can cover new combinations, but may average evidence across cases whose
difficulty or learning dynamics differ.  Interactions can express residual
dependencies while increasing sparsity.  The fixed probe protocol exposes this
trade-off by showing all-probe loss beside coverage, cold-start, and
development-supported subsets.  No policy can improve its headline metric by
discarding the items it fails to annotate.

\paragraph{The simulation is a mechanism check.}
The independent oracle removes the circularity of generating responses from
the candidate KCs under evaluation.  It still favours representations capable
of recovering its particular factorised structural regularities.  A policy
that transfers well in this world has demonstrated information reuse under a
known mechanism, not correspondence to human grammar acquisition.  Robustness
to seeds and alternative declared oracles would test engineering dependence;
human learner sequences are required for educational validity.

\paragraph{Path to tutoring dialogue.}
TSCCv2 motivates the project but is not an evaluation dataset in the present
study.  Moving from controlled items to dialogue requires three new layers:
identifying which turns are genuine assessment opportunities, assigning
context-sensitive correctness, and allowing an utterance to instantiate
multiple grammatical phenomena.  These labels should be created and audited
without using downstream KT scores to decide the ontology.  The controlled
benchmark can then serve as a unit-tested bridge, while real dialogue supplies
the independent learner evidence the current simulation lacks.

The immediate completion sequence is therefore empirical rather than
editorial: restore the verified consult-only source, run and validate the
reference condition, freeze its pre-selection outputs, execute the five
remaining policy variants through KT, complete the manual audits, and populate
Table~\ref{tab:results-contract}.  Only then should the discussion be revised
to answer the three RQs with estimates rather than expectations.
